\subsection{Условия постоянства и монотонности функций на промежутке. Локальные экстремумы}

\subsubsection*{1. Условия постоянства и монотонности}

Рассматриваются условия, связывающие знак производной функции на интервале с её монотонностью или постоянством. Основным инструментом доказательства здесь является формула Лагранжа о конечном приращении.

\begin{theorem}[Условия монотонности]
Пусть функция $f(x)$ дифференцируема на интервале $(a, b)$. Тогда выполняются следующие свойства:
\begin{enumerate}
    \item $f'(x) > 0$ на $(a, b) \implies f(x)$ строго возрастает на $(a, b)$. В обратную сторону: если $f(x)$ строго возрастает, то $f'(x) \ge 0$.
    \item $f'(x) \ge 0$ на $(a, b) \iff f(x)$ возрастает (не убывает) на $(a, b)$.
    \item $f'(x) \equiv 0$ на $(a, b) \iff f(x) \equiv \const$ на $(a, b)$.
    \item $f'(x) \le 0$ на $(a, b) \iff f(x)$ убывает (не возрастает) на $(a, b)$.
    \item $f'(x) < 0$ на $(a, b) \implies f(x)$ строго убывает на $(a, b)$. В обратную сторону: если $f(x)$ строго убывает, то $f'(x) \le 0$.
\end{enumerate}
\end{theorem}

\begin{tcolorbox}[title=Доказательство, breakable]
Доказательство состоит из двух частей: обоснование переходов слева направо (от знака производной к функции) и справа налево (от свойств функции к знаку производной).

\textbf{1. Прямые переходы (использование формулы Лагранжа):}
Рассмотрим любые две точки $x_1, x_2 \in (a, b)$, такие что $x_1 < x_2$. По теореме Лагранжа существует точка $\xi \in (x_1, x_2)$, для которой:
\[ f(x_2) - f(x_1) = f'(\xi)(x_2 - x_1) \]
Так как $x_2 - x_1 > 0$, знак приращения функции $f(x_2) - f(x_1)$ полностью определяется знаком производной $f'(\xi)$.
\begin{itemize}
    \item Если $f'(x) > 0$ на всем интервале, то $f'(\xi) > 0 \implies f(x_2) - f(x_1) > 0$, что означает строгое возрастание.
    \item Если $f'(x) \equiv 0$, то $f(x_2) - f(x_1) = 0 \implies f(x_1) = f(x_2)$, то есть функция постоянна.
\end{itemize}
Остальные пункты доказываются аналогично.

\textbf{2. Обратные переходы (использование определения производной):}
По определению производной:
\[ f'(x) = \lim_{h \to 0} \frac{f(x + h) - f(x)}{h} \]
Пусть функция $f(x)$ возрастает. Тогда при $h > 0$ имеем $f(x + h) \ge f(x)$, а при $h < 0$ имеем $f(x + h) \le f(x)$. В обоих случаях разностное отношение $\frac{f(x + h) - f(x)}{h} \ge 0$. По теореме о предельном переходе в неравенствах:
\[ f'(x) \ge 0 \]
Даже для строго возрастающей функции предел может быть равен нулю (например, для $y = x^3$ в точке $x = 0$).
\end{tcolorbox}

\subsubsection*{2. Локальные экстремумы}

\begin{definition}
Точка $x_0 \in (a, b)$ называется точкой \textbf{локального максимума} (минимума), если существует такая окрестность $U(x_0)$, что для всех $x \in U(x_0)$ выполняется неравенство $f(x) \le f(x_0)$ ($f(x) \ge f(x_0)$). Если в проколотой окрестности неравенство строгое, экстремум называется \textbf{строгим}.
\end{definition}

\begin{theorem}[Необходимое условие экстремума — теорема Ферма]
Если функция $f(x)$ имеет в точке $x_0$ локальный экстремум, то либо её производная $f'(x_0) = 0$, либо $f'(x_0)$ не существует.
\end{theorem}

\begin{theorem}[I достаточное условие экстремума]
Пусть $f(x)$ непрерывна в критической точке $x_0$ и дифференцируема в её проколотой окрестности. Если при переходе через $x_0$ производная $f'(x)$ меняет знак:
\begin{enumerate}
    \item С «$+$» на «$-$» — в $x_0$ локальный максимум.
    \item С «$-$» на «$+$» — в $x_0$ локальный минимум.
    \item Если знак производной не меняется — экстремума нет.
\end{enumerate}
\end{theorem}

\begin{theorem}[II достаточное условие (по высшим производным)]
Пусть функция $f(x)$ определена в окрестности точки $x_0$ и дифференцируема в этой точке $n$ раз ($n \ge 2$), причем:
\[ f'(x_0) = f''(x_0) = \dots = f^{(n-1)}(x_0) = 0, \quad f^{(n)}(x_0) \neq 0 \]
Тогда:
\begin{enumerate}
    \item Если $n$ — нечетное число, то в точке $x_0$ экстремума нет.
    \item Если $n$ — четное число, то в точке $x_0$ имеется локальный экстремум:
    \begin{itemize}
        \item при $f^{(n)}(x_0) > 0$ — локальный минимум;
        \item при $f^{(n)}(x_0) < 0$ — локальный максимум.
    \end{itemize}
\end{enumerate}
\end{theorem}

\begin{tcolorbox}[title=Доказательство, breakable]
Для доказательства используем формулу Тейлора с остаточным членом в форме Пеано в точке $x_0$:
\[ f(x) = f(x_0) + \sum_{k=1}^{n-1} \frac{f^{(k)}(x_0)}{k!}(x - x_0)^k + \frac{f^{(n)}(x_0)}{n!}(x - x_0)^n + o((x - x_0)^n) \]
По условию теоремы все производные до порядка $n-1$ равны нулю. Тогда приращение функции:
\[ f(x) - f(x_0) = \frac{f^{(n)}(x_0)}{n!}(x - x_0)^n + \alpha(x)(x - x_0)^n = (x - x_0)^n \left[ \frac{f^{(n)}(x_0)}{n!} + \alpha(x) \right] \]
где $\alpha(x) \to 0$ при $x \to x_0$. 

Так как $f^{(n)}(x_0) \neq 0$, то при достаточно малых $|x - x_0|$ выражение в квадратных скобках сохраняет знак производной $f^{(n)}(x_0)$.
\begin{enumerate}
    \item Если $n$ нечетно, то множитель $(x - x_0)^n$ меняет знак при переходе через $x_0$. Следовательно, разность $f(x) - f(x_0)$ также меняет знак, что по определению означает отсутствие экстремума.
    \item Если $n$ четно, то $(x - x_0)^n > 0$ в проколотой окрестности. Знак разности $f(x) - f(x_0)$ совпадает со знаком $f^{(n)}(x_0)$.
    \begin{itemize}
        \item Если $f^{(n)}(x_0) > 0$, то $f(x) > f(x_0)$ (минимум).
        \item Если $f^{(n)}(x_0) < 0$, то $f(x) < f(x_0)$ (максимум).
    \end{itemize}
\end{enumerate}
\end{tcolorbox}